\section{Poisson clocks and interface additivity}\label{sec:interfaces}

The clock construction must preserve the discrete scheduler, not merely the
set of possible moves.  Fix an initial state $S$, put $M=\max S$, and place
an independent rate-one Poisson clock at each site in the finite accessible
set $\{1,\ldots,M\}$.  A ring at an occupied site triggers
\eqref{eq:update}; a ring at an empty site is ignored, and the root has no
clock.  Every effective ring time is a stopping time for this finite clock
vector.  The strong Markov property gives fresh independent rate-one
exponential residuals there, including at a newly occupied site whose clock
may have rung earlier while empty.  If the current state has $k$ nonroot
piles, the next effective ring is consequently uniform among their $k$
sites, with total rate $k$.  Hence the embedded effective-jump chain is
exactly the chain of \cref{sec:model}.

For $0\leq a\leq b$, define $h(a,b)$ by
\begin{align}
 h(a,a)&=0,                                             \label{eq:hdiag}\\
 h(0,b)&=b,                                             \label{eq:hroot}\\
 h(a,b)&=\frac12+
 \frac{h(a-1,b)+h(a,b-1)}2 \quad(0<a<b).               \label{eq:hrec}
\end{align}
This triangular recursion uniquely defines nonnegative rational values.

\begin{theorem}[interface-additive update count]\label{thm:additivity}
Let
$S=\{0=s_0<s_1<\cdots<s_r\}$.  Then
\begin{equation}
 \boxed{\displaystyle
 \E[T_S]=\sum_{i=1}^{r}h(s_{i-1},s_i).}                \label{eq:additivity}
\end{equation}
The empty sum for $S=\{0\}$ equals zero.
\end{theorem}

\begin{proof}
Label the initial piles $0,1,\ldots,r$ from left to right and follow each
label through the site-clock arrows.  Labels that meet share one path
thereafter.  Rootward nearest-neighbor paths preserve order: they can
coalesce but cannot cross.

More precisely, at every time each current pile carries a nonempty interval
of initial labels, and the spatial order of the piles is the order of these
intervals.  This holds initially for the singleton intervals.  A move to an
empty predecessor carries one interval without changing the order.  At a
collision the occupied predecessor is the immediately preceding pile, so
the two label intervals are consecutive and their union is again an
interval.  Induction over effective rings proves the claim.  In particular,
an outside label joining either member of a designated adjacent pair does
not alter that path's future site arrows.

For $1\leq i\leq r$, let $\tau_i$ be the first meeting time of the paths
started at $s_{i-1}$ and $s_i$.  At time $t$, consecutive initial labels
belong to the same current pile exactly when their interface has closed.
The current piles are therefore the consecutive label blocks separated by
open interfaces.  One block contains the root, so the number $N_t$ of
nonroot piles satisfies the pathwise identity
\begin{equation}
 N_t=\sum_{i=1}^{r}\1_{\{\tau_i>t\}}.                  \label{eq:interfaces}
\end{equation}

Let $J_t$ count effective updates by continuous time $t$.  Conditional on
the past just before $t$, its predictable intensity is $N_{t-}$.  Thus, for
each finite $t$,
\[
 J_t-\int_0^tN_{u-}\,du
\]
is the compensated jump-count martingale and
$\E J_t=\E\int_0^tN_{u-}\,du$.  By \cref{prop:pgf}, there are at most
$\PhiPot(S)$ effective waits; each has positive finite rate before
absorption, so $J_\infty=T_S\leq\PhiPot(S)$ almost surely.  Monotone
convergence on both sides, followed by Tonelli's theorem, gives
\begin{align}
 \E T_S
 &=\E J_\infty
  =\E\int_0^\infty N_{t-}\,dt
  =\E\int_0^\infty N_t\,dt \notag\\
 &=\sum_{i=1}^{r}\E\int_0^\infty
       \1_{\{\tau_i>t\}}\,dt
  =\sum_{i=1}^{r}\E\tau_i.                             \label{eq:compensator}
\end{align}
The two time integrals agree because the jump times are Lebesgue-null.

It remains to identify one interface marginal.  Before two paths at
$0<a<b$ meet, they occupy distinct sites and read independent rate-one
clocks.  The next event takes mean time $1/2$; either the lower path moves to
$a-1$ or the upper path moves to $b-1$, each with probability $1/2$.
The strong Markov property supplies fresh exponentials even when a site was
visited earlier.  At equality the interface is closed, while a lower path at
zero is fixed and the upper path needs $b$ mean-one moves.  First-event
conditioning is exactly \eqref{eq:hdiag}--\eqref{eq:hrec}, so
$\E\tau_i=h(s_{i-1},s_i)$.  Substitution into
\eqref{eq:compensator} proves \eqref{eq:additivity}.
\end{proof}

\begin{remark}
The interface lifetimes are generally dependent.  Equation
\eqref{eq:additivity} uses only the pathwise count
\eqref{eq:interfaces}, Tonelli's theorem, and their individual means.
\end{remark}
